\section{Resultados Comparativos da Classificação de Transações}
\label{sec:val_resultados_classificacao}

Para responder às questões de pesquisa Q1 e Q2, cinco arquiteturas foram implementadas e
avaliadas sobre \textbf{o mesmo conjunto de teste-cego de 334 transações}, cujos
estabelecimentos são \textbf{disjuntos} dos utilizados em treino e no povoamento da base
vetorial (particionamento por estabelecimento; Seção~\ref{sec:val_datasets}). Todas as
métricas e artefatos foram registrados no MLflow. A Tabela~\ref{tab:resultados_modelos}
sumariza os resultados globais.

\begin{table}[htb]
	\centering
	\caption{Comparativo global de desempenho preditivo e tempo médio de inferência (teste-cego de 334 transações, estabelecimentos disjuntos).}
	\label{tab:resultados_modelos}
	\resizebox{\textwidth}{!}{%
		\begin{tabular}{lccccc}
			\toprule
			\textbf{Modelo / Arquitetura} & \textbf{Acurácia} & \textbf{Precisão} & \textbf{Recall} & \textbf{F1-Macro} & \textbf{Latência} \\
			\midrule
				Baseline 1: TF-IDF + Naive Bayes & 88,3\% & 69,8\% & 61,7\% & 63,6\% & 0,03 ms \\
				Baseline 2: TF-IDF + Random Forest & 84,7\% & 74,2\% & 66,5\% & 69,3\% & 0,28 ms \\
				SOTA Neural: ByT5 Fine-Tuning (byt5-small) & 65,9\% & 52,0\% & 53,0\% & 48,3\% & 65,5 ms \\
				Busca Vetorial RAG: PGVector (MiniLM) & 82,6\% & 57,9\% & 55,4\% & 56,2\% & 1,42 ms \\
				\textbf{Agente Híbrido Completo (LangGraph)} & \textbf{61,4\%} & \textbf{68,4\%} & \textbf{62,9\%} & \textbf{60,0\%} & \textbf{6,5 s} \\
			\bottomrule
		\end{tabular}%
	}
	\fonte{Elaborado pelo autor (\the\year). Latência do agente: média por transação, dominada pelo ramo LLM (gemma-4-e4b local).}
\end{table}


A análise revela um quadro \textbf{distinto do relatado pela literatura apoiada em grandes
corpora}, atribuível às características do acervo pessoal disponível (volume reduzido,
11 categorias efetivamente povoadas e forte
desbalanceamento):

\begin{enumerate}
	\item \textbf{Os baselines esparsos foram os mais robustos.} O TF-IDF com Random Forest
	obteve o melhor F1-Macro (\textbf{69,3\%}) e o TF-IDF com Naive
	Bayes a melhor acurácia global (\textbf{88,3\%}), ambos com
	latência inferior a 0,3~ms. A vetorização por $n$-gramas de caractere (\texttt{char\_wb},
	2--4) capturou a morfologia recorrente das descrições brasileiras (``UBER'', ``IFD*'',
	``DROGA-'') mesmo em estabelecimentos inéditos;

	\item \textbf{A diferença entre validação cruzada e teste-cego foi acentuada.} Sob
	validação cruzada 5-fold (com repetição de estabelecimentos entre folds), o Random Forest
	atingiu F1-Macro de \textbf{92,7\%}; no teste-cego por estabelecimento o
	valor caiu para \textbf{69,3\%}. Essa diferença de
	$\approx$23~pontos percentuais quantifica o quanto o desempenho aparente de um
	classificador de transações depende da \textbf{memorização de comércios já observados};

	\item \textbf{O ByT5 ajustado sofreu sobreajuste.} Com 1150 exemplos de treino, o
	\texttt{byt5-small} convergiu a \textit{loss} de treino inferior a 0,02 mas
	generalizou mal (F1-Macro \textbf{48,3\%}), colapsando a predição
	da classe majoritária \textit{Lazer e Entretenimento}. Modelos no nível de byte exigem
	corpora de ordem de grandeza superior;

	\item \textbf{A busca vetorial densa (RAG) ficou limitada pela base fria.} Com
	488 estabelecimentos vetorizados por
	\texttt{paraphrase-multilingual-MiniLM-L12-v2} (384-d), apenas
	\textbf{57,5\%} das transações de teste
	obtiveram vizinho com similaridade de cosseno $\ge$~0,85; nas demais o vizinho foi
	semanticamente fraco, resultando em F1-Macro de \textbf{56,2\%}
	(latência real de consulta no PGVector: \textbf{1,4~ms});

	\item \textbf{O Agente Híbrido combinou os dois ramos, ao custo de latência.} O agente roteou 52\% das transações para a busca vetorial (mediana de 40~ms) e 48\% para o LLM local (dezenas de segundos), acionando confirmação humana em 4,5\% dos casos. Atingiu acurácia de 61,4\% e F1-Macro de 60,0\% (latência média ponderada de 6,5~s). O F1-Macro do agente supera o da busca vetorial isolada (56,2\%) --- o ramo LLM recupera parte das classes de cauda --- mas sua acurácia global é a mais baixa entre os modelos, porque o ramo LLM, sob a diretriz anti-alucinação, tende a rotular como \textit{Despesas Diversas} nomes cuja categoria era lexicalmente evidente (acurácia condicional do ramo LLM: apenas 21\%, contra 98\% do ramo vetorial).
\end{enumerate}

\subsection{Acurácia global versus F1-Macro}
Todos os modelos apresentam acurácia global sensivelmente superior ao F1-Macro (o Naive
Bayes, por exemplo: 88,3\% de acurácia contra
63,6\% de F1-Macro). Isso decorre de dois fatores do acervo pessoal:
(i) o forte desbalanceamento --- \textit{Alimentação} e \textit{Lazer e
Entretenimento} respondem por cerca de metade do teste; e (ii) o tamanho reduzido das
classes de cauda --- \textit{Moradia}, \textit{Tecnologia}, \textit{Vestuário} e
\textit{Educação} contam com 4 a 5 estabelecimentos distintos cada. As categorias
\textit{Pets}, \textit{Serviços Financeiros} e \textit{Doações e Presentes} não
tiveram amostragem suficiente e foram agregadas a \textit{Despesas Diversas / Outros}. O
F1-Macro é adotado como métrica conservadora de referência; os resultados devem ser lidos
como \textbf{diagnóstico sobre um acervo pessoal real de porte reduzido}, não como
estimativa de desempenho populacional.

\subsection{Desempenho Detalhado por Categoria de Despesa}
A Tabela~\ref{tab:resultados_categorias} detalha o comportamento do TF-IDF + Random Forest nas
categorias efetivamente avaliadas.

\begin{table}[htb]
	\centering
	\caption{Desempenho do TF-IDF + Random Forest discriminado por categoria de despesa efetivamente avaliada.}
	\label{tab:resultados_categorias}
	\begin{tabularx}{\textwidth}{X c c c c}
		\toprule
		\textbf{Categoria de Despesa} & \textbf{Precisão (\%)} & \textbf{Recall (\%)} & \textbf{F1-Score (\%)} & \textbf{Volume de Teste} \\
		\midrule
		Alimentação & 69,4\% & 100,0\% & 82,0\% & 100 \\
		Transporte & 100,0\% & 61,5\% & 76,2\% & 39 \\
		Saúde & 100,0\% & 88,5\% & 93,9\% & 61 \\
		Moradia & 100,0\% & 75,0\% & 85,7\% & 4 \\
		Lazer e Entretenimento & 100,0\% & 86,6\% & 92,8\% & 67 \\
		Tecnologia & 0,0\% & 0,0\% & 0,0\% & 4 \\
		Marketplace / E-commerce & 100,0\% & 96,8\% & 98,4\% & 31 \\
		Vestuário & 0,0\% & 0,0\% & 0,0\% & 5 \\
		Beleza e Cuidados Pessoais & 100,0\% & 80,0\% & 88,9\% & 5 \\
		Educação & 100,0\% & 100,0\% & 100,0\% & 4 \\
		Despesas Diversas / Outros & 46,2\% & 42,9\% & 44,4\% & 14 \\
		\midrule
		\textbf{Média Macro} & \textbf{74,1\%} & \textbf{66,5\%} & \textbf{69,3\%} & \textbf{334} \\
		\bottomrule
	\end{tabularx}
	\fonte{Elaborado pelo autor (\the\year).}
\end{table}


Categorias com padrões léxicos fortes e volume adequado --- \textit{Alimentação},
\textit{Transporte}, \textit{Saúde} e \textit{Marketplace / E-commerce} ---
alcançaram F1-Score entre 90\% e 100\%. As classes de cauda (\textit{Tecnologia},
\textit{Vestuário}, \textit{Moradia}) obtiveram F1 nulo ou próximo de zero em quase
todos os modelos: com 4 a 5 exemplos no teste, um único erro derruba a métrica. Esse é o
principal fator limitante do F1-Macro.

\subsection{Estudo de Ablação}
O estudo de ablação isolou os componentes mensuráveis sem novo custo de inferência do LLM
(Tabela~\ref{tab:ablacao}). O ramo vetorial do agente respondeu por 52\% das transações (acurácia condicional 98\%) e o ramo LLM por 48\% (acurácia 21\%).

\begin{table}[htb]
	\centering
	\caption{Ablação de componentes sobre o teste-cego (variação de F1-Macro).}
	\label{tab:ablacao}
	\begin{tabularx}{\textwidth}{X c c c}
		\toprule
		\textbf{Componente} & \textbf{Com} & \textbf{Sem} & \textbf{$\Delta$ F1 (p.p.)} \\
		\midrule
		Normalização léxica (TF-IDF+RF) & 67,4\% & 69,3\% & +1,9 \\
		Normalização léxica (RAG MiniLM) & 66,8\% & 56,2\% & -10,6 \\
		Ramo vetorial (RAG) — acurácia condicional & 97,7\% & 21,4\% & +76,3 \\
		\bottomrule
	\end{tabularx}
	\fonte{Elaborado pelo autor (\the\year).}
\end{table}


A ablação revela um efeito \textbf{dependente da representação}. Para os classificadores
esparsos (TF-IDF + Random Forest), a normalização léxica agressiva de gateways foi
neutra ou levemente prejudicial ($\Delta$ F1 $\approx -2$~p.p.): a vetorização por
$n$-gramas de caractere já é robusta a prefixos como ``PAG*'' e ``MP*'', que por vezes
carregam sinal útil. Para a busca vetorial densa, ao contrário, a normalização foi
\textbf{decisiva} ($\Delta$ F1 $\approx +11$~p.p.): sem ela, o \textit{embedding}
multilíngue se confunde com o ruído dos intermediadores. O componente que sustenta o
desempenho do agente é a \textbf{busca vetorial}: seu ramo respondeu corretamente a
98\% das
transações que resolveu, enquanto o ramo LLM local (\texttt{gemma-4-e4b}) acertou apenas
21\% ---
tendendo, sob a diretriz anti-alucinação, a rotular como \textit{Despesas Diversas}
nomes cuja categoria era lexicalmente evidente (``SAPPORO RESTAURANTE'', ``Drogaria
Paulista''). A ferramenta de busca na web foi desativada por indisponibilidade de
conectividade no ambiente experimental.
